% Introduction draft: fast photometric/geometric convergence is the thesis.
% The Carve paragraph and the third contribution are intentionally modular.

\section{Introduction}
\label{sec:intro}
% chaehyeon: 아래 첫 두문단은 요약해서 한문단으로 써보세요. 
3D Gaussian Splatting (3DGS) has established itself as a highly effective representation for photorealistic novel view synthesis~\cite{kerbl20233dgs}. By modeling scenes with explicit, colored Gaussian primitives and rendering them via a tile-based differentiable rasterizer, it provides a compelling combination of high visual fidelity and real-time rendering speeds. Consequently, 3DGS is increasingly utilized across a diverse range of visual applications. 
To achieve this level of quality, the representation requires careful optimization of numerous parameters—including position, scale, rotation, color, and opacity—for millions of primitives. Since this optimization process is computationally intensive and relies on dense multi-view constraints, 3DGS frameworks are predominantly trained offline using completely collected datasets.

However, autonomous robots and wearable augmented-reality systems (e.g., Aria glasses) are fundamentally bound to the present moment. For these active agents, high reconstruction quality achieved only after an exhaustive offline refinement stage is inadequate; they intrinsically require online mapping to construct and reliably refine a map as new observations continuously arrive. 
Adapting 3DGS to this online constraint exposes a critical vulnerability. While recent SLAM systems strive to process streams quickly, "real-time" operation is an inherently ambiguous metric that fluctuates heavily depending on the available GPU hardware. Therefore, rather than engineering a system merely to meet an arbitrary processing speed, the true objective shifts to achieving \emph{fast photometric and geometric convergence} for practical field deployments. Regardless of the specific hardware environment, the system is forced to make a newly observed region structurally and visually reliable as quickly as possible under whatever computational budget is given.
% ------------ 여기까지가 한문단

Building a reliable online Gaussian map inherently requires two complementary objectives: establishing a geometrically sound starting point and rapidly optimizing it under strict computational limits. 
Recent advancements have approached these goals from various angles. For instance, EDGS~\cite{edgs} demonstrates that dense image correspondences can provide a robust geometric initialization, laying the groundwork for high-quality reconstruction. At the optimizer level, 3DGS-LM introduces second-order solvers to accelerate convergence~\cite{hollein2025gslm}, while systems like CaRtGS and PGSR incorporate allocation strategies and geometric consistencies to maximize the efficiency of available compute~\cite{cartgs, chen2024pgsr}. 
While these approaches successfully highlight the necessity of learning efficiency and structural reliability, relying solely on advanced optimizers or dense initialization is insufficient. To truly pursue fast convergence in a streaming environment, the fundamental question shifts from \emph{how} to optimize to \emph{when} and \emph{which} observations should enter the training pipeline.

%------------------- 두번째 문단 끝 -----------------



The rate of map convergence relies critically on the view selection policy. Current online systems, such as Photo-SLAM~\cite{huang2024photoslam} and VIGS-SLAM~\cite{vigsslam}, heavily depend on frontend-selected keyframes for map optimization. While bounding computation to local keyframe windows (as seen in MonoGS and HI-SLAM2~\cite{matsuki2024monogs,zhang2025hislam2}) efficiently reuses selected observations, discarding the intermediate frames leaves valuable texture and viewpoint information unexploited. 
Conversely, naively injecting all available frames into the optimization pool dilutes the limited update budget, leaving each view insufficiently refined. Uniform sampling from an ever-growing pool fails to compensate for the varying entry times of new observations, systematically delaying the convergence of recently added regions. Hence, achieving rapid photometric improvement requires decoupling the mapping-view set from the tracking trajectory. By dynamically coupling the inclusion of new observations to actual completed optimization work and meticulously balancing training opportunities across views, a limited computational budget can be strictly directed to bring the appearance of newly observed regions to high quality sooner.

A strong geometric starting point is essential, but it does not guarantee that structures remain accurate during subsequent color-driven optimization. Due to the shape-radiance ambiguity of alpha compositing, misplaced Gaussians can effortlessly replicate target colors by forming semi-transparent clouds in empty space~\cite{deng2022dsnerf, stablegs, wang2025efags}. One might assume that geometric losses can resolve this by simply translating (i.e., updating the positions of) these floating artifacts back to the correct surface. However, relying on positional gradients exposes a fundamental flaw in the 3DGS architecture. Because the rasterizer depends on strict depth-sorting~\cite{kerbl20233dgs}, shifting a Gaussian along the camera ray frequently alters its sorting order relative to its neighbors, causing abrupt, discontinuous spikes in the transmittance integral. Furthermore, it is inherently ambiguous whether a floater should be pushed toward the foreground object or pulled into the background. Attempting to spatially navigate this highly non-convex landscape wastes critical online updates. Therefore, rather than struggling to move misplaced Gaussians, an online method should directly control their opacity, explicitly suppressing density in observed free space to rapidly eliminate spurious structures without relying on unstable positional gradients~\cite{yu2024gof, guedon2023sugar, pagas2026}.

To achieve this fast convergence within a bounded online budget, we build upon a DROID-SLAM-based frontend~\cite{teed2021droid}. While recent works have explored dense geometric initialization~\cite{edgs}, efficient second-order solvers~\cite{hollein2025gslm}, and compute-aware optimization allocations~\cite{cartgs, chen2024pgsr}, translating these foundations into continuous map improvement requires strict control over which observations enter training. Therefore, we tightly couple our mapping policy to the actual optimization progress. For photometric refinement, we introduce an Optimization-Guided View Growth strategy that expands the mapping-view set strictly at the rate of completed updates. This is complemented by our Entropy-Regularized Count Balancing (ERCB) sampler, which dynamically corrects severe training deficits among views while preserving the diversity of sampling without replacement. For geometric refinement, we apply a causal free-space Carve loss that directly acts on opacity. By continuously penalizing opacity where depth evidence dictates empty space, we enforce structural consistency immediately after a region is observed.

We validate whether this unified architecture fundamentally accelerates map convergence rather than merely offering a fast rendering speed. Through controlled incremental evaluations, we compare our rendering quality and structural accuracy against recent state-of-the-art online systems, notably VIGS-SLAM~\cite{vigsslam}. To strictly isolate the efficacy of our proposed mapping policies from the inherent benefits of our dense initialization, we keep the initialization protocol completely frozen across all ablation studies. This rigorous experimental design explicitly demonstrates how our view scheduling and opacity-directed constraints extract superior, highly converged map quality from the exact same computational budget. Our contributions are summarized as follows:
\begin{itemize}
    \item Building upon a dense, frontend-driven initialization, we formulate an online 3DGS mapping framework optimized strictly for \emph{fast photometric and geometric convergence} under limited compute.
    \item We propose an Optimization-Guided View Growth policy that paces the inclusion of causal frames according to completed mapping work.
    \item We introduce Entropy-Regularized Count Balancing (ERCB) to correct severe optimization imbalances while preserving the essential diversity of sampling without replacement.
    \item We design a causal, depth-supported Carve loss that targets opacity directly, rapidly suppressing free-space artifacts and bypassing the severe gradient discontinuities associated with positional optimization.
\end{itemize}

% GEOMETRY-MODULE SWITCH:
% If fixed opacity or another controller wins the final ablation, replace only
% the geometry sentences in paragraphs 5--6 and contribution 3. The problem
% statement, photometric path, and convergence-based evaluation remain valid.